\documentclass[11pt]{article}


\usepackage[inline]{asymptote}
\usepackage{asypictureB}
%\usepackage[default]{frcursive}
\usepackage[a4paper,text={17.7cm,25.8cm},centering,top=21mm]{geometry} 
\usepackage{amsfonts}
%\usepackage{amsmath}
\usepackage{amssymb,amsmath,stackrel,wasysym,mathrsfs}
\usepackage{multicol}
\usepackage{pdflscape}% girar una pagina, hay que usar \begin{landscape}En este punto inicia la hoja en vertical \end{landscape}
%\usepackage[mathcal]{euscript}
%\usepackage{amssymb}
\usepackage{esvect}%para la flecha de vector 

\usepackage[T1]{fontenc}
\usepackage[spanish,es-noshorthands]{babel}
%\usepackage[latin1,ansinew]{inputenc}\usepackage{calligra}
%\usepackage[pdftex]{graphicx}
%\usepackage{etex}
\usepackage{calligra}
\usepackage{array}
 \usepackage{booktabs}
 \usepackage{ifsym}
\usepackage{pgf,tikz,pgfplots}
\usepackage{tkz-graph}%paquete (di)grafos
\usepackage{tkz-euclide}% Para realizar diagramas que coloreen conjuntos
\usetikzlibrary[decorations.text]% Decorar curvas
%\usetikzlibrary{matrix}
\usetikzlibrary{arrows,matrix}
%\usetkzobj{all}

\pgfplotsset{compat=1.15}
\usepackage{mathrsfs}
\usetikzlibrary{arrows}
\usepackage{makecell} %Paquete para controlar altura de celdas
\usepackage{verbatim} 
\usetikzlibrary{babel}%para no general error al usar babel y tikz simultaneamente.
\usetikzlibrary{patterns,patterns.meta}%%%%%%%%%%%para  pintar de forma rrallada un rectangulo

\usepackage{pictex}
\usepackage{verbatim}%me sirve para usar comentarios con mucha lineas 
\usepackage{amstext}
\usepackage{enumerate}
%\usepackage{enumitem}
\usepackage{amsthm}
\usepackage{xcolor}
\usepackage{mathdots}% para tener mas opciones de puntos suspensivos

%\usepackage{pgf,tikz}%,pgfplots
%\pgfplotsset{compat=1.15}
\usepackage{mathrsfs}
%\usetikzlibrary{arrows}
%\usepackage{etex,tkz-euclide} 
%\usetkzobj{all}
\usepackage{color}

\usepackage{cancel}
\newcommand{\Cancelar}[2][black]{{\color{#1}\cancel{\color{black}#2}}}%%% para cancelar con colores recordar usar mayusculas
\usepackage{lscape}%paquete para girar hojas











\newcommand{\dis}{\displaystyle}
\newcommand{\R}{\mathbb{R}}
\newcommand{\re}{\mathbb{R}}
\newcommand{\I}{\mathbb{I}}
\newcommand{\Q}{\mathbb{Q}}
\newcommand{\Z}{\mathbb{Z}}
\newcommand{\C}{\mathbb{C}}
\newcommand{\N}{\mathbb{N}}
\newcommand{\K}{\mathbb{K}}
%\newcommand{\B}{\mathcal{B}}
\newcommand{\V}{\mathbb{V}}
\newcommand{\can}{\mathcal{C}}
\newcommand{\M}{\mathbb{M}}
\newcommand{\Pol}{\mathcal{P}}
\newcommand{\moduloprop}{\hookrightarrow_{\hspace{-2.5mm}\mbox{\tiny$^{_\neq}$}}}

\newtheorem{ejer}{Ejercicio}[]

\newtheorem{teo}{Teorema}%[section]
\newtheorem{lema}[teo]{Lema}
\newtheorem{propo}[teo]{Proposición}
\newtheorem{defi}[teo]{Definición}
\newtheorem{ejem}[teo]{}
\newtheorem*{ejemsinn}{}
\newtheorem{obs}[teo]{Observación}%[section]
\newtheorem{acla}{Aclaración}%[section]
\newtheorem{act}{Actividad}%%%%%%%%%%%%%%%%
\newtheorem*{defsinn}{{\bf Definición}:}
\newtheorem*{ejemplosinn}{{\bf Ejemplo}:}
\newtheorem*{proposinn}{{\bf\sc Proposición}}
\newtheorem*{obssinn}{Observación:}
\newtheorem*{corosinn}{Corolario}
\newtheorem*{propsinn}{Propiedades:}
%\newenvironment{proof} { \noindent {\bf Demostración: }\rm}{$\quad \hfill \blacksquare $}
\usepackage{setspace} % (es una alternativa a \renewcommand{\baselinestretch}{1.3}) usando \doublespacing \onehalfspace \singlespace \spacing{1.5}  se va cambiando segun uno quiera
\spacing{1.3}

%\newenvironment{proof}[Demostración:]
\newenvironment{sol}{ \noindent {\bf Solución:} \rm}{%\hfill \protect\tikz[overlay,yshift=-4pt]\protect\node[xshift=8mm, yshift=13mm, blue]() {\huge $\checkmark$ };
}

%%%%%%%%%%%%%%%%%%%%%%%%%%%%%%%%%%%%%%%%


%\renewcommand{\baselinestretch}{1.5}
%%%%%%%%%%%%%%%%%%%%%%%%%%%%%%%%%%%%%%%%%%%%%%%%%%%%%%%%%%%%%%%%%%%%%%%%%%%%%%%%%%%%%%%%%%%%%%%%%%%%%%%%%%%%%%%%%%%%%%%%%%%%%%%%%%%%%%%%%%%%%%%%%%%%%%%%%%%%%%%%%%%%%%%%%%%%%%%%%%%%%%%%%%%%%%%%%%%%%%%%%%%%%%%%%%%%%%%%%%%%%%%%%%%%%%%%%%%%%%%%%%%%%%%%%%%%%%%%%%%%%%%%%%%%%%%%%%%%%%%%
\usepackage[framemethod=TikZ]{mdframed}%Paquete para estilos de los bordes de los ambientes Importante este paquete (mdframed) define un entorno  que trata automáticamente los saltos de página en el texto enmarcado.
            \mdfsetup{skipabove=1\baselineskip, skipbelow=1.5\baselineskip ,             frametitlefont= \sffamily\bfseries , needspace=4\baselineskip , splittopskip=1.5\baselineskip} %
            \mdfsetup{apptotikzsetting={\tikzset{mdfbackground/.append style={draw=none}}}} 
%%%%%%%%%%%%%%%%%%%%%%%%%%%%%%%%%%%%%%%%%%%%%%%%%%%%%%%%%%%%%%
%
%Estilo para Cuentas auxiliares
%
%%%%%%%%%%%%%%%%%%%%%%%%%%%%%%%%%%%%%%%%%%%%%%%%%%%%%%%%%%%%%%
\mdfdefinestyle{Cuentaauxi}{%skipabove=10pt, skipbelow=0,
userdefinedwidth=1.02\textwidth,align=center,
  everyline=true,
  ignorelastdescenders=true,
  middlelinewidth=1pt,
  linecolor=black!20,
  outerlinewidth=1pt,
  %leftline=false,topline=false
  %rightline=false,bottomline=false,
  %backgroundcolor=black!20,
  innerleftmargin=2mm, innerrightmargin=2mm,
  innerbottommargin=3mm,
  innertopmargin=1mm,
  frametitleaboveskip=2mm,
  frametitlebelowskip=2mm,
  %frametitlerule=false,
  %repeatframetitle=false,
 % outerlinewidth=6pt,outerlinecolor=blue!20,
 pstricksappsetting={\addtopsstyle{mdfouterlinestyle}{linestyle=dashed}},
 %firstextra={\node[xshift=5cm,right,draw=White, line width=1.5pt,%star,star points=10,star point ratio=1.2, 
 %minimum size=15mm, fill=white] at (P-|O) {\tikzpicture\draw[draw=black!20,decoration=snake, fill=white] (0,0) --(2.5,0) decorate{--(2.5,1)}--(0,1)decorate{--cycle};%\includegraphics[width=10mm]{.pdf}
 %  \endtikzpicture
 %}; },
 %singleextra={\node[xshift=5cm,right,draw=White, line width=1.5pt,%star,star points=10,star point ratio=1.2, minimum size=15mm,
  %fill=white] at (P-|O) {\tikzpicture\draw[draw=black!20,decoration=snake, fill=white] (0,0) --(2.5,0) decorate{--(2.5,1)}--(0,1)decorate{--cycle};%\includegraphics[width=10mm]{.pdf}    
 %  \endtikzpicture
 %}; }
}
%%%%%%%%%%%%%%%%%%%%%%%%%%%%%%%%%%%%%%%%%%%%%%%%%%%%%%%%%%%%%%
%
%Estilo para Citar
%
%%%%%%%%%%%%%%%%%%%%%%%%%%%%%%%%%%%%%%%%%%%%%%%%%%%%%%%%%%%%%%
\mdfdefinestyle{citar}{
userdefinedwidth=.83\textwidth ,align=center,
skipabove=5pt, skipbelow=5pt,
  everyline=true,
  ignorelastdescenders=true,
  middlelinewidth=1pt,
  linecolor=blue!60,
  outerlinewidth=2pt,
  %leftline=false,
  topline=false,
  rightline=false,
  bottomline=false,
  %backgroundcolor=black!20,
  innerleftmargin=5mm, innerrightmargin=5mm,
  innerbottommargin=3mm,
  innertopmargin=1mm,
  leftmargin= 14mm,
  %frametitleaboveskip=2mm,
  %frametitlebelowskip=2mm,
  %frametitlerule=false,
  %repeatframetitle=false,
 % outerlinewidth=6pt,outerlinecolor=blue!20,
 pstricksappsetting={\addtopsstyle{mdfouterlinestyle}{linestyle=dashed}},
 %firstextra={\node[xshift=5cm,right,draw=White, line width=1.5pt,%star,star points=10,star point ratio=1.2, 
 %minimum size=15mm, fill=white] at (P-|O) {\tikzpicture\draw[draw=black!20,decoration=snake, fill=white] (0,0) --(2.5,0) decorate{--(2.5,1)}--(0,1)decorate{--cycle};%\includegraphics[width=10mm]{.pdf}
 %  \endtikzpicture
 %}; },
 %singleextra={\node[xshift=5cm,right,draw=White, line width=1.5pt,%star,star points=10,star point ratio=1.2, minimum size=15mm,
  %fill=white] at (P-|O) {\tikzpicture\draw[draw=black!20,decoration=snake, fill=white] (0,0) --(2.5,0) decorate{--(2.5,1)}--(0,1)decorate{--cycle};%\includegraphics[width=10mm]{.pdf}    
 %  \endtikzpicture
 %}; }
}
%%%%%%%%%%%%%%%%%%%%%%%%%%%%%%%%%%%%%%%%%%%%%%%%%%%%%%%%%%%%%%%%%%%%%%%%%%%%%%%%%%%%%%%%%%%%%%%%%%%%%%%%%%%%%%%%%%%%%%%%%%%%%%%%%%%%%%%%%%%%%%%%%%%%%%%%%%%%%%%%%%%%%%%%%%%%%%%%%%%%%%%%
%%%%%%%%%%%%%%%%%%%%%%%%%%%%%%%%%%%%%%%%%%%%%%%%%%%%%%%%%%%%%%%%%%%%%%%%%%%%%%%%%%%%%%%%%%%%%%%%%%%%%%%%%%%%%%%%%%%%%%%%%%%%%%%%%%%%%%%%%%%%%%%%%%%%%%%%%%%%%%%%%%%%%%%%%%%%%%%%%%%%%%%%%%%%%%%%%%%%%%%%%%%%%%%%%%%%%%%%%%%%%%%%%%%%%%
%%%%%%%%%%%%%%%%%%%%%%%%%%%%%%%%%%%%%%%%%%%%%%%%%%%%%%%%%%%%%%%%%%%%%%%%%%%%%%%%%%%%%%%%%%%%%%%%%%%%%%%%%%%%%%%%%%%%%%%%%%%%%%%%%%%%%%%%%%%%%%%%%%%%%%%%%%%%%%%%%%%%%%%%%%%%%%%%%%%%%%%%%%%%%%%%%%%%%%%%%%
\newcommand{\Dem}{\protect\tikz[overlay,yshift=-4pt]\protect\draw[line width=1.5pt,line cap=round,color=gray] (0,0)[out=14,in=180] to node [shift=(90:6.pt),color=black] {\large\bf Dem:} (7:1); \qquad \quad }
\DeclareRobustCommand{\michi}{{\mathpalette\irchi\relax}}
\newcommand{\irchi}[1]{\raisebox{\depth}{\large$#1\chi$}} % inner command, used by \rchi



%%%%%%%%%%%%%%%%%%%%%%%%%%%%%%%%%%%%%%%%%%%%%%%%%%%%%%%%%%%%%%%%%%%%%%%%%%%%%%%%%%%%%%%%%%%%%%%%%%%%%%%%%%%%%%%%%%%%%%%%%%%%%%%%%%%%%%%%%%%%%%%%%%%%%%%%%%%%%%%%%%%%%%%%%%%%%%%%%%%%%%%%%%%%%%%%%%%%%%%%%%%%%%%%%%%%%%%%%%%%%%%%%%%%%%%%%%%%%%%%%%%%%%%%%%%%%%%%%%%%%%%%%%%%%%%%%%%%%%%%%%%%%%%%%%%%%%%%%%%%%%%%%%%%%%%%%%%%%%%%%%%%%%%%%%%%%%%%%%%%%%%%%%%%%%%%%%%%%%%%%%%%%%%%%%%%%%%%%%%%%%%%%%%%%%%%%%%%%%%%%%%%%%%%%%%%%%%%%%%%%%%%%%%%%%%%%%%%%%%%%%%%%%%%%%%%%%%%%%%%%%%%%%%%%%%%%%%%%%%%%%%%%%%%%%%%%%%%%%%%%%%%%%%%%%%%
% comando aclaraciones  

\newcommand{\aclaracion}[2]{

\begin{center}
\definecolor{gris}{rgb}{0.5,0.5,0.5}
\begin{tikzpicture}[line cap=round,line join=round,scale=1]
%\draw [color=gris,dash pattern=on 1pt off 1pt, xstep=1cm,ystep=1cm] (0.1,0.1) grid (6.3,6.3);
%\draw[color=black] (0.,0.) -- (\linewidth,0.);
\node at (.23,{.9*#2}) [color=black,above right] {\bf #1};
\foreach \y in {1,2,...,#2}
\draw[shift={(0,{.9*\y})},color=gris,opacity=.5] (0.,0.) -- ({.99*\linewidth},0.);
\end{tikzpicture} 
\end{center}
}





%%%%%%%%%%%%%%%%%%%%%%%%%%%%%%%%%%%%%%%%%%%%%%%%%%%%%%%%%%%%%%%%%%%%%%%%%%%%%%%%%%%%%%%%%%%%%%%%%%%%%%%%%%%%%%%%%%%%%%%%%%%%%%%%%%%%%%%%%%%%%%%%%%%%%%%%%%%%%%%%%%%%%%%%%%%%%%%%%%%%%%%%%%%%%%%%%%%%%%%%%%%%%%%%%%%%%%%%%%%%%%%%%%%%%%%%%%%%%%%%%%%%%%%%%%%%%%%%%%%%%%%%%%%%%%%%%%%%%%%%%%%%%%%%%%%%%%%%%%%%%%%%%%%%%%%%%%%%%%%%%%%%%%%%%%%%%%%%%%%%%%%%%%%%%%%%%%%%%%%%%%%%%%%%%%%%%%%%%%%%%%%%%%%%%%%%%%%%%%%%%%%%%%%%%%%%%%%%%%%%%%%%%%%%%%%%%%%%%%%%%%%%%%%%%%%%%%%%%%%%%%%%%%%%%%%%%%%%%%%%%%%%%%%%%%%%%%%%%%%%%%%%%%%%%%%%
% comando respuesta 

\newcommand{\res}[1]{

\begin{center}
\definecolor{gris}{rgb}{0.5,0.5,0.5}
\begin{tikzpicture}[line cap=round,line join=round,scale=1]
%\draw [color=gris,dash pattern=on 1pt off 1pt, xstep=1cm,ystep=1cm] (0.1,0.1) grid (6.3,6.3);
%\draw[color=black] (0.,0.) -- (\linewidth,0.);
\node at (.23,{.9*#1}) [color=black,above right] {\bf respuesta:};
\foreach \y in {1,2,...,#1}
\draw[shift={(0,{.9*\y})},color=gris,opacity=.5] (0.,0.) -- ({.99*\linewidth},0.);
\end{tikzpicture} 
\end{center}
}

%%%%%%%%%%%%%%%%%%%%%%%%%%%%%%%%%%%%%%%%%%%%%%%%%%%%%%%%%%%%%%%%%%%%%%%%%%%%%%%%%%%%%%%%%%%%%%%%%%%%%%%%%%%%%%%%%%%%%%%%%%%%%%%%%%%%%%%%%%%%%%%%%%%%%%%%%%%%%%%%%%%%%%%%%%%%%%%%%%%%%%%%%%%%%%%%%%%%%%%%%%%%%%%%%%%%%%%%%%%%%%%%%%%%%%%%%%%%%%%%%%%%%%%%%%%%%%%%%%%%%%%%%%%%%%%%%%%%%%%%%%%%%%%%%%%%%%%%%%%%%%%%%%%%%%%%%%%%%%%%%%%%%%%%%%%%%%%%%%%%%%%%%%%%%%%%%%%%%%%%%%%%%%%%%%%%%%%%%%%%%%%%%%%%%%%%%%%%%%%%%%%%%%%%%%%%%%%%%%%%%%%%%%%%%%%%%%%%%%%%%%%%%%%%%%%%%%%%%%%%%%%%%%%%%%%%%%%%%%%%%%%%%%%%%%%%%%%%%%%%%%%%%%%%%%%%
% comando ejemplo para completar 

\newcommand{\ejparacompletar}[1]{

\begin{center}
\definecolor{gris}{rgb}{0.5,0.5,0.5}
\begin{tikzpicture}[line cap=round,line join=round,scale=1]
%\draw [color=gris,dash pattern=on 1pt off 1pt, xstep=1cm,ystep=1cm] (0.1,0.1) grid (6.3,6.3);
%\draw[color=black] (0.,0.) -- (\linewidth,0.);
\node at (.23,{.9*#1}) [color=black,above right] {\bf Ejemplo:};
\foreach \y in {1,2,...,#1}
\draw[shift={(0,{.9*\y})},color=gris,opacity=.5] (0.,0.) -- ({.99*\linewidth},0.);
\end{tikzpicture} 
\end{center}
}


%%%%%%%%%%%%%%%%%%%%%%%%%%%%%%%%%%%%%%%%%%%%%%%%%%%%%%%%%%%%%%%%%%%%%%%%%%%%%%%%%%%%%%%%%%%%%%%%%%%%%%%%%%%%%%%%%%%%%%%%%%%%%%%%%%%%%%%%%%%%%%%%%%%%%%%%%%%%%%%%%%%%%%%%%%%%%%%%%%%%%%%%%%%%%%%%%%%%%%%%%%%%%%%%%%%%%%%%%%%%%%%%%%%%%%%%%%%%%%%%%%%%%%%%%%%%%%%%%%%%%%%%%%%%%%%%%%%%%%%%%%%%%%%%%%%%%%%%%%%%%%%%%%%%%%%%%%%%%%%%%%%%%%%%%%%%%%%%%%%%%%%%%%%%%%%%%%%%%%%%%%%%%%%%%%%%%%%%%%%%%%%%%%%%%%%%%%%%%%%%%%%%%%%%%%%%%%%%%%%%%%%%%%%%%%%%%%%%%%%%%%%%%%%%%%%%%%%%%%%%%%%%%%%%%%%%%%%%%%%%%%%%%%%%%%%%%%%%%%%%%%%%%%%%%%%%
% comando conclusión 

\newcommand{\conclus}[1]{

\begin{center}
\definecolor{gris}{rgb}{0.5,0.5,0.5}
\begin{tikzpicture}[line cap=round,line join=round,scale=1]
%\draw [color=gris,dash pattern=on 1pt off 1pt, xstep=1cm,ystep=1cm] (0.1,0.1) grid (6.3,6.3);
%\draw[color=black] (0.,0.) -- (\linewidth,0.);
\node at (.23,{.9*#1}) [color=black,above right] {\bf Conclusión:};
\foreach \y in {1,2,...,#1}
\draw[shift={(0,{.9*\y})},color=gris,opacity=.5] (0.,0.) -- ({.99*\linewidth},0.);
\end{tikzpicture} 
\end{center}
}

%%%%%%%%%%%%%%%%%%%%%%%%%%%%%%%%%%%%%%%%%%%%%%%%%%%%%%%%%%%%%%%%%%%%%%%%%%%%%%%%%%%%%%%%%%%%%%%%%%%%%%%%%%%%%%%%%%%%%%%%%%%%%%%%%%%%%%%%%%%%%%%%%%%%%%%%%%%%%%%%%%%%%%%%%%%%%%%%%%%%%%%%%%%%%%%%%%%%%%%%%%%%%%%%%%%%%%%%%%%%%%%%%%%%%%%%%%%%%%%%%%%%%%%%%%%%%%%%%%%%%%%%%%%%%%%%%%%%%%%%%%%%%%%%%%%%%%%%%%%%%%%%%%%%%%%%%%%%%%%%%%%%%%%%%%%%%%%%%%%%%%%%%%%%%%%%%%%%%%%%%%%%%%%%%%%%%%%%%%%%%%%%%%%%%%%%%%%%%%%%%%%%%%%%%%%%%%%%%%%%%%%%%%%%%%%%%%%%%%%%%%%%%%%%%%%%%%%%%%%%%%%%%%%%%%%%%%%%%%%%%%%%%%%%%%%%%%%%%%%%%%%%%%%%%%%%
% boton calculadora 

\newcommand{\boton}[2]{
\tikzpicture{  \node[draw,color=#1,fill=#1,text=white,,rounded corners=3mm,text width=7mm,text height =3mm,align=center,midway]  at (0,0) { \large \bf #2 };
}
\endtikzpicture
}

\usepackage[framemethod=TikZ]{mdframed}%Paquete para estilos de los bordes de los ambientes Importante este paquete (mdframed) define un entorno  que trata automáticamente los saltos de página en el texto enmarcado.
            \mdfsetup{skipabove=2\baselineskip,skipbelow=2\baselineskip,frametitlefont=\sffamily\bfseries, needspace=4\baselineskip, splittopskip=1.5\baselineskip} 
            \mdfsetup{apptotikzsetting={\tikzset{mdfbackground/.append style={draw=none}}}} 
            
            
%%%%%%%%%%%%%%%%%%%%%%%%%%%%%%%%%%%%%%%%%%%%%%%%%%%%%%%%%%%%%%
%
%Estilo  para cajas
%
%%%%%%%%%%%%%%%%%%%%%%%%%%%%%%%%%%%%%%%%%%%%%%%%%%%%%%%%%%%%%%

\mdfdefinestyle{teoSty}{backgroundcolor=blue!10, linecolor=blue, linewidth=3pt,
skipabove=4pt,
skipbelow=3pt, 
innerleftmargin=3ex, innerrightmargin=3ex, innertopmargin=3mm, innerbottommargin=3mm, innermargin=15pt, outermargin=-15pt, needspace={0.2\textheight},roundcorner=10pt}

\mdfdefinestyle{ObsSty}{backgroundcolor=red!10, linecolor=red, linewidth=2pt,
skipabove=4pt,
skipbelow=3pt, 
innerleftmargin=3ex, innerrightmargin=3ex, innertopmargin=3mm, innerbottommargin=3mm, innermargin=5pt, outermargin=-5pt, needspace={0.2\textheight},roundcorner=10pt}

\mdfdefinestyle{EjerSty}{backgroundcolor=blue!5!green!5!, linecolor=blue!50!green!50!, linewidth=2pt,
skipabove=4pt,
skipbelow=3pt, 
innerleftmargin=3ex, innerrightmargin=3ex, innertopmargin=3mm, innerbottommargin=3mm, innermargin=5pt, outermargin=-5pt, needspace={0.2\textheight},roundcorner=10pt}



%%%%%%%%%%%%%%%%%%%%%%%%%%%%%%%%%%%%%%%%%%%%%%%%%%%%%%%%%%%%%%
%
%Estilo  para lineas de relleno usando el paquete patterns y la libreria  patterns.meta
%
%%%%%%%%%%%%%%%%%%%%%%%%%%%%%%%%%%%%%%%%%%%%%%%%%%%%%%%%%%%%%%

\tikzdeclarepattern{
  name=mylines,
  parameters={
      \pgfkeysvalueof{/pgf/pattern keys/size},
      \pgfkeysvalueof{/pgf/pattern keys/angle},
      \pgfkeysvalueof{/pgf/pattern keys/line width},
  },
  bounding box={
    (0,-0.5*\pgfkeysvalueof{/pgf/pattern keys/line width}) and
    (\pgfkeysvalueof{/pgf/pattern keys/size},
0.5*\pgfkeysvalueof{/pgf/pattern keys/line width})},
  tile size={(\pgfkeysvalueof{/pgf/pattern keys/size},
\pgfkeysvalueof{/pgf/pattern keys/size})},
  tile transformation={rotate=\pgfkeysvalueof{/pgf/pattern keys/angle}},
  defaults={
    size/.initial=5pt,
    angle/.initial=45,
    line width/.initial=.4pt,
  },
  code={
      \draw [line width=\pgfkeysvalueof{/pgf/pattern keys/line width}]
        (0,0) -- (\pgfkeysvalueof{/pgf/pattern keys/size},0);
  },
}





\newcommand{\semejante}[1]{\stackrel[#1]{}{\longleftrightarrow}}% \thicksim para decir la Semejanza de matrices la variable #1  indica la operacion elemental 
\newcommand{\OptipoUno}[1]{\stackrel[#1]{}{=}}



\newcommand{\indicador}[4]{
{\tikz[remember picture,overlay]\coordinate (#1) at #2;}{\tikz[remember picture,overlay]\coordinate (#3) at #4;} 
}

\newcommand{\marka}[2]{
{\tikz[remember picture,overlay]\coordinate (#1) at #2;}
}




\newcommand{\cof}{ \mbox{cof}}

\newcommand{\paralela}{\rotatebox[origin=c]{-10}{\Large$\parallel$}}

\usepackage{fancyhdr}
\fancypagestyle{miestilo}{
\fancyhead[L]{}
\fancyhead[C]{}
\fancyhead[R]{}
\fancyfoot[L]{}
\fancyfoot[C]{}
\fancyfoot[R]{pág \thepage}
\renewcommand{\headrulewidth}{0pt}
\renewcommand{\footrulewidth}{0pt}
}
%%%%%%%%%%%%
%%%%%%%%%%%%
\renewcommand{\footrulewidth}{1pt}
\pagestyle{fancy}\markboth{}{}
\renewcommand{\headrulewidth}{.05pt}
\renewcommand{\headrule}{{\color{gray}%
\hrule width\headwidth height\headrulewidth \vskip-\headrulewidth}}
\renewcommand{\footrule}{{\color{gray}%
\vskip-\footruleskip\vskip-\footrulewidth
\hrule width\headwidth height\footrulewidth\vskip\footruleskip}}
\renewcommand{\footrulewidth}{.50pt}
\fancyhead[L]{\color{gray}{\footnotesize \sc  \materia} }
\fancyhead[C]{Ingeniería en Petróleo}
\fancyhead[R]{\color{gray}{\sc \footnotesize  Universidad Nacional Del Comahue}}
\fancyfoot[R]{pág. \thepage}
\fancyfoot[C]{}
\setlength{\headheight}{15pt}

%\usepackage{wrapfig}

\newcommand{\co}{\mathbb{C}}
%\newcommand{\re}{\mathbb{R}}
\newcommand{\q}{\mathbb{Q}}
\newcommand{\z}{\mathbb{Z}}
\newcommand{\n}{\mathbb{N}}
\newcommand{\ve}{\mathbb{V}}
\newcommand{\w}{\mathbb{W}}
\newcommand{\be}{{\cal B}}
\newcommand{\ca}{{\cal C}}
\newcommand{\pqx}{\mathbb{Q}[x]}
\newcommand{\prx}{\mathbb{R}[x]}
\newcommand{\pcx}{\mathbb{C}[x]}
%\newcommand{\izg}{\big\langle}
%\newcommand{\rangle}{\big\rangle}
%%%%%%%%%%%%%%%%%%%%%%%%%%%%%5

%%%%%%%%%%%%%%%%%%%%%%%%%%%%%%%%%%%%5
\newcommand{\materia}{\sc Álgebra y Geometría I}%separadas con '&'
%%%%%%%%%%%%%%%%%%%%%%%%%%%%%%%%%%%%%%%
% se define el encabezado, 
% el parametro es el texto que va despues de: trabajo practico $N^\circ$... (agregar numero y titulo)
%
%%%%%%%%%%%%%%%%%%%%%%%%%%%%%%%%%\newcommand{\encabezado}[1]{%
\fancyfoot[L]{\vspace*{-3.5mm}\color{gray}{\small T.P. 3: Sistema de Ecuaciones Lineales }}
\thispagestyle{miestilo} 
\newcommand{\encabezado}[1]{%
\fancyfoot[L]{\vspace*{-3.5mm}\color{gray}{\small T. P. #1}}
%\thispagestyle{miestilo}%
%--------------------logodpto //texto // logouni------------------
%--------------------logodpto------------------

\vspace*{-18mm}\hspace*{-11mm} \begin{minipage}{2.7cm}
\begin{center} 
     \includegraphics[width=2.7cm, height=2.7cm]{logodpto}
\end{center}  
\end{minipage}\hspace*{-15mm} \begin{minipage}{15.9cm}
\begin{center}
\begin{tabular}{c @{ - } c}
 \materia \\
\carreras
\end{tabular}

\cuatri

\medskip

TRABAJO PRÁCTICO N$^\circ$#1 

\end{center}

\end{minipage}\hspace{-14mm} \begin{minipage}{2.5cm}
  \begin{center} 
    \includegraphics[width=2.5cm, height=2.5cm]{logouni}
  \end{center}  
\end{minipage}
%-----------------------------------------------------

%----------------------------------------------------------------
\hrule
%---------------------------------------------------------------
}%%



\pagestyle{fancy}
\def\identacion{-2.5mm}%mueve la identacion de la enumeracion de los ejercicios.


\usepackage{asypictureB}
\begin{asyheader}
settings.outformat="png";
settings.render=8;
import graph3;
import smoothcontour3;

pen bpen = rgb(0.75, 0.07, 0.01);
//material m = material(diffusepen=0.7bpen, ambientpen=bpen+opacity(0.8), emissivepen=0.3*bpen,specularpen=0.999white, shininess=1.0);      
\end{asyheader}


\usepackage{multicol}
\setlength{\columnsep}{7mm}




\begin{document}

\setlength{\headheight}{15pt}

\vspace*{-2.5cm}
%--------------------logodpto------------------
\begin{figure}[htbp]
  \begin{minipage}[r][][t]{0.15\textwidth}
\begin{center} 
 % Universidad Nacional del Comahue.  
 \includegraphics[width=2.2cm, height=2.2cm]{imagen/inglogo.jpg} \ \ 
\vspace{1.8cm}
\end{center}  
\end{minipage}\quad
% texto central
\begin{minipage}[r][][t]{0.35\textwidth}
\begin{center}
\vspace*{-25mm}
{\bf Universidad Nacional
del Comahue} \\
Facultad de Ingeniería
\end{center}
\end{minipage}\quad\begin{minipage}[r][][t]{0.27\textwidth}
\begin{center}
\vspace*{-25mm}
{\bf  Álgebra y Geometría I}
Ingeniería en Petróleo\\
Primer Cuatrimestre 2026
\end{center}
\end{minipage}\quad
%----------------------logouni-------------------------
\begin{minipage}[r][][t]{0.15\textwidth}
  \begin{center} 
     \ \   \includegraphics[width=2.2cm, height=2.2cm]{imagen/Logo_UNco.jpg}
     \vspace{1.8cm}
  \end{center}  
\end{minipage}
%-----------------------------------------------------
\end{figure}
\vspace*{-30 mm}

\begin{center}
{\bf \Large{ Trabajo Práctico 3: Sistemas de ecuaciones lineales. }}
\end{center}
\hrule

 


\begin{enumerate}[\hspace{-7mm} 1) ]  
\item  Para cada una de las siguientes ecuaciones, dar la solución general y dos soluciones particulares.

\begin{enumerate}
\begin{multicols}{3}
    \item $2x - 4y = -6$
    \item $3x - y + 2z = 5$
    \item $-x_1+3x_2-\frac{1}{4}x_3-2x_4+x_5=2$
\end{multicols}
\end{enumerate}





\item Dados los siguientes sistemas de ecuaciones lineales, 

\vspace{-8mm}
{\singlespace
    \begin{enumerate}[i)]
    \begin{multicols}{2}
        \item  $\left\{\begin{array}{l}       y = 2x - 3 \\ 
        x + 4y = 6 \\ 
        3x - 7 = -y \\
        -\tfrac{1}{2}x - \tfrac{1}{2}y + \tfrac{3}{2} = 0
        \end{array}\right.$

\medskip 
        
        \item $\begin{bmatrix} -2 & 1 \\ 2 & 3 \\ 2 & -1 \end{bmatrix} \cdot \begin{bmatrix} x \\ y \end{bmatrix} = \begin{bmatrix} -3 \\ 2 \\ -4 \end{bmatrix}$


\columnbreak %Cambio de columna 

        
        \item  $\begin{bmatrix} -1 \\ \tfrac{1}{2} \end{bmatrix}x 
        +
        \begin{bmatrix}3 \\ -\tfrac{3}{2} \end{bmatrix} y
        = \begin{bmatrix} -24 \\ 12 \end{bmatrix}$

\bigskip  
        
        \item $\left\{ \begin{array}{rl}
        y &= -2x \\
        x + 3y &= 0 \\
        2y - 2x &= 1
        \end{array}\right.$   
        
        \end{multicols}
    \end{enumerate} 
 }   
    
    \begin{enumerate}[a)]
    \item Resolver analíticamente y gráficamente. Relacionar la compatibilidad del sistema con la posición relativa de las rectas en el plano.

    
    \item  Determinar si los pares ordenados siguientes son soluciones de los sistemas indicados:
    \begin{itemize}
    \begin{multicols}{3}
        \item $(-1, -1)$ para el sistema (I).
        \item $(6, -6)$ para el sistema (III).
        \item $(0, 0)$ para el sistema (IV).
        
    \end{multicols} 
    \end{itemize}
\end{enumerate}















\item Resolver mediante el método de Gauss o de Gauss-Jordan y clasificar los siguientes sistemas de ecuaciones. Escribir en cada caso el conjunto solución. Si el sistema es compatible indeterminado, dar además de la solución general, dos soluciones particulares.

\vspace{-8mm}
{\singlespace
\begin{enumerate}[i)]
\begin{multicols}{2}
    \item $\left\{\begin{array}{rl}
    2x + 3y - z &= 4 \\
    x - 2y + z &= -7 \\
    -3x + 2y + 2z &= 14 \\
    3x + y &= -3
    \end{array} \right.$

    
    \item $\begin{bmatrix} 1 & 1 & -1 \\ 2 & 4 & -1 \\ -1 & 1 & 2 \end{bmatrix} \cdot \begin{bmatrix} x \\ y \\ z \end{bmatrix} = \begin{bmatrix} 0 \\ 0 \\ 0 \end{bmatrix} $

    \columnbreak % Forzar cambio de columna si es necesario


    \item  $
    \begin{bmatrix} 1 \\ -2 \end{bmatrix} x + \begin{bmatrix} -2 \\ 4 \end{bmatrix} y + \begin{bmatrix} 7 \\ -14 \end{bmatrix} z = \begin{bmatrix} -1 \\ -2 \end{bmatrix}
    $

    
 \medskip    
    
    
    \item $\left\{ \begin{array}{rl}   x + y - z &= 0 \\
    4y + 2x - z &= 0 \\
    2z - x + y &= 0
    \end{array} \right.$




\medskip   


    
    
    \item $\begin{bmatrix} 1 \\ 3 \end{bmatrix} x + \begin{bmatrix} 2 \\ -1 \end{bmatrix} y + \begin{bmatrix} 3 \\ -5 \end{bmatrix} z  + \begin{bmatrix} -4 \\ -5 \end{bmatrix} w= \begin{bmatrix} 5 \\ 1  \end{bmatrix}
    $

\end{multicols}

\end{enumerate}

}






 \item Resolver los siguientes sistemas, si es posible, aplicando la inversibilidad de matrices. En caso de no ser posible, resolverlos de otra manera.
\begin{enumerate}[i) ]
\begin{multicols}{2}
    \item $\begin{bmatrix} 1 \\ 7 \end{bmatrix} x + \begin{bmatrix} -1 \\ 5 \end{bmatrix} y + \begin{bmatrix} 5 \\ -13 \end{bmatrix} z = \begin{bmatrix} -2 \\ 10 \end{bmatrix}
    $


\medskip
    
    \item  $ \left\{\begin{array}{rl} 
    3x + y - 2z &= 10 \\
    x + z + 2y &= 5 \\
    2z - x &= -3
    \end{array}\right. $


  \columnbreak % Forzar cambio de columna si es necesario

    \item $\begin{bmatrix}\tfrac{1}{5} & -1 & 1 & -2\\
0 & 1 & -1 & 2\\
-1 & 5 & -5 & 10\\
-1 & 7 & -7 & 14\end{bmatrix} \cdot 
\begin{bmatrix} x\\y\\z\\w \end{bmatrix}
= 
\begin{bmatrix}
    4 \\-1 \\ -20 \\-22
\end{bmatrix}
$ 
\, \\    
\end{multicols}
    
\end{enumerate}














\item Dado el sistema $\left\{ \begin{matrix}
x - y + 3z &= 2 \\
x - 3y + 3z &= 6 \\
-x + 5y - 3z &= -10
\end{matrix} \right.$, determinar si los siguientes conjuntos son o no soluciones correspondientes al sistema dado.

\begin{enumerate}
\begin{multicols}{2}
    \item $S = \left\{ \left(-3z, -2, z\right) \, ,\, z \in \mathbb{R}\right\}$
    \item $S = \left\{\left(-3z, -2, -\tfrac{1}{3}x  \rule{0mm}{4mm}\right) \, ,\, x, z \in \mathbb{R}\right\}$
    \item $S = \left\{\left( x, -2, -\tfrac{1}{3}x \rule{0mm}{4mm}\right) \, ,\, x \in \mathbb{R}\right\}$
    \item[d)] $S = \left\{(-3z, y, z),\, y, z \in \mathbb{R}\right\}$
\end{multicols}
\end{enumerate}


















\item  En los siguientes incisos se da la forma escalonada reducida de un sistema de ecuaciones lineales. Determinar la compatibilidad del sistema utilizando el Teorema de Rouché-Fröbenius y, en caso de ser compatible, dar el conjunto solución poniéndole por nombre a las variables \(x_1, x_2, \ldots, x_n\) según corresponda. 


\vspace{-8mm}
{\singlespace
\begin{enumerate}[a) ]
\begin{multicols}{3}
    \item  $\left[
\begin{array}{ccc|c}
    1 & 0 & 0 & 0 \\
    0 & 1 & 0 & 0 \\
    0 & 0 & 0 & 1
    \end{array}\right] $

\bigskip
    
    \item  $\left[
\begin{array}{ccc|c}
    1 & 0 & -2 & 5 \\
    0 & 1 & 3 & -1 \\
    0 & 0 & 0 & 0 \\
    0 & 0 & 0 & 0
    \end{array}\right] $

\columnbreak %Cambio de columna 

    
    \item $\left[
\begin{array}{ccc|c}
    1 & 0 & 0 & -8 \\
    0 & 1 & 0 & 0 \\
    0 & 0 & 1 & 7
    \end{array}\right]$

    
    \item  $\left[
\begin{array}{cccc|c}
    1 & 0 & 0 & 0 & 0 \\
    0 & 1 & 0 & 2 & 2 \\
    0 & 0 & 1 & 0 & 0
    \end{array}\right]$ 
    

\columnbreak %Cambio de columna 

   
    \item  $\left[
\begin{array}{cccc|c}
    1 & 1 & 1 & 0 & 1 \\
    0 & 0 & 0 & 1 & -1 \\
    0 & 0 & 0 & 0 & 0
    \end{array}\right]$
    
    \item $\left[
\begin{array}{ccc|c}
    1 & 0 & 0 & 4 \\
    0 & 1 & 0 & -3 \\
    0 & 0 & 1 & 0 \\
    0 & 0 & 0 & 0
    \end{array}\right]$
\end{multicols}
\end{enumerate}
}
























\item Determinar para qué valores de \(k \in \mathbb{R}\) cada uno de los siguientes sistemas es: %discutir los sistemas de ecuaciones.
\begin{itemize}
\begin{minipage}{.36\textwidth}
       \item Compatible determinado (CD),
\end{minipage}
\begin{minipage}{.36\textwidth}
    \item Compatible indeterminado (CI),

\end{minipage}\begin{minipage}{.17\textwidth}
    \item Incompatible (I).
\end{minipage}
\end{itemize}


\begin{enumerate}[a) ]
\begin{multicols}{2}
    \item  $\left\{
    \begin{array}{rl}
    5x + 3y + 4z &= 3 \\
    5y + 2z &= 4 \\
    (k^2 - 9)z &= k + 3
    \end{array}\right.$ 
    
    
    \medskip
    
    \item  $
    \left[\begin{smallmatrix}
    1 \\[1mm] 2 \\[1mm]  1
    \end{smallmatrix}\right] x +
    \left[\begin{smallmatrix}
    1 \\[1mm]  3 \\[1mm]  k
    \end{smallmatrix}\right] y +
    \left[\begin{smallmatrix}
    -1 \\[1mm]  2 \\[1mm]  3
    \end{smallmatrix}\right] z =
    \left[\begin{smallmatrix}
    1 \\[1mm]  3 \\[1mm]  2
    \end{smallmatrix}\right]$

\medskip

    \item $
    \begin{bmatrix}
    1 & 2 & -3 \\
    3 & -1 & 5 \\
    4 & 1 & -(k^2 - 6)
    \end{bmatrix}
    \cdot
    \begin{bmatrix}
    x \\ y \\ z
    \end{bmatrix} =
    \begin{bmatrix}
    4 \\ 2 \\ k + 4
    \end{bmatrix}$

 \medskip
 
    \item  $\left\{ 
    \begin{array}{rl}
    x \quad + \qquad 2y \quad + \qquad 3z \quad \, &= 1 \\
    2x + (k + 4)y + (k^2 + k + 6)z &= k + 3
    \end{array}\right.$ 

\columnbreak 
    
    \item $
    \left[\begin{smallmatrix}
    k \\[1mm]  1 \\[1mm]  1
    \end{smallmatrix}\right] x +
    \left[\begin{smallmatrix}
    1 \\[1mm]  k \\[1mm]  1
    \end{smallmatrix}\right] y +
    \left[\begin{smallmatrix}
    1 \\[1mm]  1 \\[1mm]  1
    \end{smallmatrix}\right] z =
    \left[\begin{smallmatrix}
    0 \\[1mm]  0 \\[1mm]  0
    \end{smallmatrix}\right]$  


\bigskip 


     \item  $\left\{
    \begin{array}{rl}
    x + y + z + (k - 1)w &= 0\\
    -x \quad - \quad ky \quad -\quad z &= -2k \\
    kx + 2y + 2z + k^2 w &= 2
    \end{array}\right.$ 

\bigskip


    \item $
    \begin{bmatrix}
    1 & 3 & -2 \\
    0 & k & 0 \\
    0 & k &  k^2 - 25 \\
    2 & 6+k & k^2-29 
    \end{bmatrix}
    \cdot
    \begin{bmatrix}
    x \\ y \\ z
    \end{bmatrix} =
    \begin{bmatrix}
    5 \\ 0 \\ k + 5 \\ k+15
    \end{bmatrix}$
\end{multicols}
\end{enumerate}











\item La siguiente matriz es la matriz de coeficientes asociada a un sistema de ecuaciones lineales homogéneo:
\[
A =
\begin{bmatrix}
2 & 2 & 1 & 0 \\
0 & 3 & 0 & 3k \\
0 & 0 & 1 & -k \\
0 & 0 & 0 & k^2 - 25
\end{bmatrix}
\]

\begin{enumerate}
    \item[a)] Determinar para qué valores de \(k\) el sistema es compatible determinado, compatible indeterminado o incompatible.
    \item[b)] Considerando \(k = 5\), escribir la solución del sistema de ecuaciones. En caso de ser compatible indeterminado, escribir dos soluciones particulares.
\end{enumerate}

\end{enumerate}

\newpage
\section*{Ejercicios para Examen Final}


\begin{enumerate}[\hspace{-7mm} 1) ]
\item \begin{enumerate}
		\item \begin{enumerate}
			\item[i)] Resolver el sistema \(
			\left \{
			\begin{array}{l}
				x+y-\tfrac{1}{2}z=5 \\
				x-3y-z=6
			\end{array}
			\right .
			\) 
			\item[ii)] Añadir una tercer ecuación para que el sistema resulte:\ ii)1) C.D.\   ii)2) C.I.\  ii)3) Inc.
		\end{enumerate} 
	
	\item Dado un sistema de dos ecuaciones lineales con tres incógnitas, analizar si siempre es posible agregar una ecuación para que resulte C.D. Justificar
\end{enumerate}


\item Determinar qué valores deben tomar $a$ y $b$ para que:

\begin{enumerate}
	\item El sistema \(\left[\begin{matrix}
		1 & b\\ 
		b&a^2
	\end{matrix}\right]
	\cdot
	\left[\begin{matrix}
		x\\
		y
	\end{matrix}\right]
	=
	\left[\begin{matrix}
		0\\0
			\end{matrix}
	\right]\) tenga infinitas soluciones. Dar dos ejemplos de valores posibles para $a$ y $b$.
	
	\item \((-1;2;1)\) sea solución del sistema \(\left[\begin{matrix}
		-1\\1\\2b
	\end{matrix}\right]
	x + 
	\left[\begin{matrix}
		a\\-1\\0
	\end{matrix}\right]
	y + 
	\left[\begin{matrix}
		1\\2\\a-b
	\end{matrix}\right]
	z =
	\left[\begin{matrix}
		-1\\-1\\2
	\end{matrix}
	\right]\)
	
	\item El conjunto \(A=\left\{(x;y;z)=\left(k;\tfrac{11}{4}-\tfrac{5}{4}k;\tfrac{9}{4}+\tfrac{1}{4}k\right),k,\in\mathbb{R}\right\}\) sea solución del sistema \(
	\left \{
	\begin{array}{l}
		5x+4y=a \\
		x+y+z=b
	\end{array}
	\right .
	\) 
	
\end{enumerate}
\item Dar, para cada uno de los siguientes incisos y en caso de ser posible, un sistema de ecuaciones lineales que cumpla con las condiciones pedidas. Justificar.
\begin{enumerate}
	\item  Un sistema homogéneo de dos ecuaciones con dos incógnitas que sea compatible  indeterminado.
	\item Un sistema homogéneo de tres ecuaciones con tres incógnitas que sea compatible indeterminado.
	\item Un sistema de dos ecuaciones con dos incógnitas que sea incompatible.
	\item Un sistema de dos ecuaciones con tres incógnitas que sea incompatible.
	\item Un sistema de dos ecuaciones con tres incógnitas que sea compatible determinado.
\end{enumerate}

\item Justificar que el sistema $\left\{
    \begin{array}{rl}
    5a x +y + az  &= 3 \\
    x +ay +z &= 4 \\
    bx + bz &= 0 
    \end{array}\right.$ es compatible indeterminado para cualquier $a,\, b \in \R$ y hallar una solución general para los $a\neq 0$ y $b\neq 0$


\item Dado el sistema de ecuaciones $\left\{
    \begin{array}{rl}
    x + \beta y + \alpha z &= 1 \\
    x + \alpha \beta y + z &= \beta\\
    \alpha x + \beta y + z  &= \alpha 
    \end{array}\right.$, analizar la compatibilidad del sistema si $\beta=0$.
\end{enumerate}





\newpage
\subsection*{ANEXO - Problemas De Aplicación.}
\begin{enumerate}[\hspace{-7mm} 1) ]
\item Una compañía petrolera está evaluando la perforación de un pozo. Los ingresos por la venta de petróleo se estiman en \( U\$D \,50 \) por barril, menos \( U\$D \,10,000 \) fijos. Además, los costos operativos se estiman en \( U\$D \,20 \) por barril, más \( U\$D \,5,000 \) en costos fijos. ¿Cuál debe ser la cantidad de petróleo producido para que los ingresos sean iguales a los costos? ¿De cuánto es el ingreso y el costo?


\item  En cada uno de los siguientes problemas, plantear el sistema de ecuaciones lineales, resolverlo
y responder.
\begin{enumerate}[a) ]
    \item Una refinería desea producir 1000 barriles de gasolina de alta calidad y 800 barriles de diésel de calidad premium. Cuenta con dos tipos de petróleo crudo: A (alta calidad) y B (baja calidad). Se sabe que para producir los barriles de gasolina deseados, durante
el proceso de refinación se utiliza el $70 \%$ de barriles del petróleo A y el$ 50 \%$  de barriles del petróleo B. Además, para producir los barriles de diesel, se utiliza el $30 \%$ de barriles del petróleo A y el $50 \%$ de barriles del petróleo B. ¿Cuántos barriles de cada tipo de petróleo crudo se necesitan para lograr la producción deseada? 



    \item En una red de tuberías de una instalación petrolera, el flujo de un tipo de fluido a través de una tubería es de 2 barriles por minuto más la mitad del flujo del otro tipo de fluido. El flujo del segundo tipo de fluido es de 200 barriles por minuto menos la mitad del flujo del primer tipo. ¿Cuáles son los flujos (en barriles por minuto) de cada tipo de fluido?

    
    \item Una compañía minera extrae mineral de dos minas. El extraído de la mina 1 contiene un 1\% de níquel y un 2\% de cobre, mientras que el de la mina 2 contiene un 1\% de níquel y un 5\% de cobre. Si la empresa busca obtener 4 toneladas de níquel y 9 de cobre, ¿qué cantidad de mineral deberá extraer de cada mina?
\end{enumerate}












\item 
Supongamos que se está estudiando el flujo de petróleo en un sistema de tres tuberías interconectadas y se sabe que la presión de cada tubería está relacionada con el flujo de petróleo a través de la misma:
\begin{itemize}
    \item Para el primer tramo, el doble de la presión de la tubería 1 debe ser igual a \( 100 \, \text{psi} \) más la presión de la tubería 2.
    \item En el segundo tramo, el triple de la presión de la tubería 2 debe ser igual a \( 50 \, \text{psi} \) más la presión de la tubería 1 y la tubería 3.
    \item Por último, para el tercer tramo, el doble de la presión de la tubería 3 debe ser \( 80 \, \text{psi} \) más la presión de la tubería 2.
\end{itemize}
Determinar la presión que debe tener cada una de las tres tuberías.






















\item En una refinería de petróleo, se desean mezclar tres tipos de combustibles: $A,\, B$ y $C$, para obtener una mezcla con un octanaje deseado. Se sabe que el combustible A tiene un octanaje de $85$, el combustible $B$ tiene un octanaje de $90$ y el combustible $C$ tiene un octanaje de $95$. Si se desea preparar una mezcla de $1000$lt, ¿cuáles son las cantidades de cada tipo de combustible que deben mezclarse para obtener la mezcla con un octanaje de $92$?.


{\bf Observación:} El octanaje indica la calidad y estabilidad del combustible. Es un indice de pureza (es el porcentaje multiplicado 100).









\end{enumerate}



\end{document}